\documentclass[12pt]{article}

\usepackage[linesnumbered,ruled,vlined]{algorithm2e}

\begin{document}

\section*{Pseudocode}

% ----------------------------------------
% Wrapper
% ----------------------------------------
\begin{algorithm}[H]
\caption{Wrapper}
\KwData{Shared Snapshot c}
\KwResult{Calls \texttt{DistAlgorithm} and save the invocation and
response in the \texttt{Snapshot} using \texttt{write} and \texttt{snapshot}}

\SetKwProg{Fn}{Function}{:}{}

\Fn{execute($op$, $param$)}{
    \texttt{c.write(id,op(param))}\;
    \texttt{response}=Execute $op(param)$ of \texttt{DistAlgorithm}\;
    \texttt{c.snapshot(response)}\;
    %\texttt{c.write(id,response)}\;
    Return \texttt{response}\;
}
\end{algorithm}

% ----------------------------------------
% Verifier
% ----------------------------------------
\begin{algorithm}[H]
\caption{Verifier}
\KwData{History of operations $X_E$, History in the snapshot $W_E$,
Shared snapshot \texttt{c}, 
DecisionAlgorithm \texttt{d}}
\KwResult{Obtains the current execution written in the
    snapshot and verifies if it is linearizable}

\SetKwProg{Fn}{Function}{:}{}

\Fn{execute()}{
    $W_E =$ \texttt{c.scanAll()} \;
    $X_E =$ Obtain the well-formed execution\;
    \texttt{boolean verdict} = \texttt{d.verify()}\;
    Return \texttt{boolean verdict}\;
}
\end{algorithm}



% ----------------------------------------
% Abstract class Snapshot
% ----------------------------------------
\begin{algorithm}[H]
\caption{Abstract Class Snapshot}
\KwData{Shared array of atomic registers}
\KwResult{Provides \texttt{write}, \texttt{snapshot} 
and \texttt{scanAll} operations}

\SetKwProg{Fn}{Function}{:}{}

\Fn{write($id$, $event$)}{
    Write $event$ to the position corresponding to $id$\;
}

\Fn{snapshot($response$)}{
    Read the last event written in all registers, 
    add the $response$ and return void\;
}

\Fn{scanAll()}{
    Return all the operations written in the snapshot\;
}
\end{algorithm}

% ----------------------------------------
% Implementation: CollectSnapshot DAG
% ----------------------------------------
\begin{algorithm}[H]
\caption{Implementation: CollectSnapshotDAG}
\KwData{SWMR Node with $n$ pointers}
\KwResult{DAG of invocation and response events}

\SetKwProg{Fn}{Function}{:}{}

\Fn{init($n$)}{
    Create $n$ Nodes, one for each process\;
}

\Fn{write($id$, $event$)}{
    Create Node($event$, $id$, $counter$)\;
    Add Node\;
}

\Fn{snapshot(response)}{
    Create Node($response$)\;
    Add pointers to the last node of each process \;
    Add Node\;
}

\Fn{scanAll()}{
    Read the DAG using a topological sort to obtain a list\;
}
\end{algorithm}

% ----------------------------------------
% Implementation: CollectSnapshot FandINC
% ----------------------------------------
\begin{algorithm}[H]
\caption{Implementation: CollectSnapshotFAInc}
\KwData{$n$ SWMR queues, Shared atomic counter $ts$}
\KwResult{$n$ queues with every invocation and response
event with a timestamp}

\SetKwProg{Fn}{Function}{:}{}


\Fn{init($n$)}{
    $ts$ is an atomic integer\;
    Create $n$ SWMR queues, one for each process\;
}

\Fn{write($id$, $event$)}{
    $i= ts.fetchAndInc()$\;
    $queue[id].enq(i,event)$\;
}

\Fn{snapshot(response)}{
    $i=ts.get()$\;
    $queue[id].enq(i,response)$\;
}

\Fn{scanAll()}{
    Read the Queues\;
}
\end{algorithm}

% ----------------------------------------
% Implementation: SnapshotWF (Wait-Free)
% ----------------------------------------
\begin{algorithm}[H]
\caption{Implementation: SnapshotWF}
\KwData{Shared registers with multi-version support}
\KwResult{Wait-free snapshot}

\SetKwProg{Fn}{Function}{:}{}

\Fn{write($id$, $event$)}{
    Write $event$ to the position corresponding to $id$\;
}

\Fn{snapshot($response$)}{
    Read the last event written in all registers, 
    add the $response$ and return void\;
}

\Fn{scanAll()}{
    Return all the operations written in the snapshot\;
}
\end{algorithm}


\end{document}